\documentclass[10pt,twocolumn,letterpaper]{article}
\usepackage[margin=0.70in,top=0.70in,bottom=0.74in,columnsep=0.24in]{geometry}
\usepackage{amsmath,amssymb,bm,booktabs,array,tabularx}
\usepackage[colorlinks=true,linkcolor=black,citecolor=black,urlcolor=black]{hyperref}

\setlength{\parindent}{0pt}
\setlength{\parskip}{2.8pt}
\renewcommand{\arraystretch}{1.05}

\makeatletter
\renewcommand{\section}{\@startsection{section}{1}{\z@}%
  {6.5pt \@plus 2pt \@minus 1pt}%
  {2.2pt}%
  {\normalfont\normalsize\bfseries}}
\makeatother

\begin{document}

\twocolumn[{%
\begin{center}
{\large\bfseries QUARCC Purple: Encompassing LASSO, Semivariance Risk, and Relax-and-Round Construction in a Market-Neutral US Equity Fund}\\[6pt]
{\normalsize QUARCC Purple team}\\
{\small McGill--FIAM Asset Management Hackathon 2026, Final Round $\cdot$ Code and locked results: \href{https://github.com/NicholasXydis/QuarccPurple}{github.com/NicholasXydis/QuarccPurple}}\\[3pt]
{\small November 2026}
\end{center}
\vspace{2pt}
\begin{center}
\begin{minipage}{0.91\textwidth}
\small
\textbf{Abstract.} We describe QUARCC Purple, a simulated \$1 billion US equity fund built for the final round of the McGill--FIAM Asset Management Hackathon 2026 under a market-neutral mandate benchmarked to the 3-month T-bill plus 4\% a year. Monthly walk-forward forecasts come from a value-weighted Encompassing LASSO (Han, He, Rapach, and Zhou, 2019) over 147 characteristics from Jensen, Kelly, and Pedersen (2023), augmented with a Stambaugh, Yu, and Yuan (2012) mispricing composite, a machine-learning forecast of next-quarter earnings surprises, and publication-lagged short interest. Targets are CAPM residual ranks. A horizon-matched factor model with directional specific semivariances prices risk, and a relax-and-round quadratic program forms a book twice a month that is traded over five sessions. From January 2021 to August 2026 (68 months, 1,421 sessions), after modelled costs, the fund grows \$1.000bn to \$1.467bn: a 7.00\% CAGR, 3.96\% volatility, Sharpe ratio 0.91, maximum drawdown $-2.99\%$, and S\&P 500 beta 0.012. It does \emph{not} beat its benchmark: the information ratio against T-bill $+4\%$ is $-0.11$ (standard error 0.42). The forecasts are walk-forward, but construction settings were chosen with the evaluation window visible, so these results are a simulation rather than a clean holdout or a live record.
\end{minipage}
\end{center}
\vspace{6pt}
}]

\section{Mandate, Data, and Point-in-Time Clocks}

The mandate requires 200--500 names, gross exposure at most 200\%, net exposure within $\pm 50\%$, a 5\% single-name limit, sector net exposure within $\pm 20\%$, and no short position below a \$500m market capitalisation. The benchmark is the 3-month T-bill plus 4\% a year, compounded monthly.

The panel is the competition's US equity file of 147 characteristics from Jensen, Kelly, and Pedersen (2023; hereafter JKP), with daily stock returns. Licensed data cannot be redistributed; the public repository contains the code, the locked outputs needed to verify the headline results, and small public series. Every input is stamped with the time it became public: SEC filings at their acceptance time, rolled to the next session when accepted after the close, and exchange short interest at eight trading days after settlement. Rows flagged as revised are excluded because their first-reported value is unavailable. Forward-looking split factors are not used.

At each month-end the estimation universe requires $|\mathrm{prc}_{i,t}| \ge \$5$ and 126-day average daily dollar volume of at least \$2m. Short eligibility additionally requires market equity of at least \$500m, and candidate positions require at least \$3bn. A book is formed at each month-end and at the first session on or after the 15th of the following month, giving 136 rebalances. Trading starts one session after formation and runs over five sessions.

\section{Signals and Encompassing LASSO Forecaster}

\textbf{Three added signals ($147 \to 150$ predictors).}
\begin{enumerate}
\item \textit{Mispricing composite.} Following Stambaugh, Yu, and Yuan (2012), we average the signed uniform ranks of ten anomalies available in the panel, omitting distress: O-score, share issuance, net equity issuance, operating accruals, net operating assets, asset growth, and investment (negative sign); and momentum, gross profitability, and return on assets (positive sign). A stock needs at least five of the ten.
\item \textit{Earnings-surprise forecast.} Bernard and Thomas (1989) condition on realised surprises. We instead train an annual walk-forward LightGBM regressor on the 147 characteristics to predict next quarter's as-filed standardised unexpected earnings before the announcement.
\item \textit{Short interest.} Following Boehmer, Huszar, and Jordan (2010), we use short interest over shares outstanding at the latest settlement date that is public by the formation close.
\end{enumerate}

\textbf{Standardisation and mid-month refresh.} Following Green, Hand, and Zhang (2017), each predictor is winsorised at its 1st and 99th cross-sectional percentiles and standardised, with missing values set to zero:
\begin{equation}
z_{i,j,t} = \frac{\mathrm{clip}\big(x_{i,j,t},\, Q_{0.01}(x_{\cdot,j,t}),\, Q_{0.99}(x_{\cdot,j,t})\big) - \mu_{j,t}}{\sigma_{j,t}}. \label{eq:ghz}
\end{equation}
At the mid-month formation $T_{\text{mid}}$ we recompute nine 21-day price signals (one-month return, volatility, idiosyncratic volatility, maximum daily returns, and skewness measures) from daily returns and substitute them, holding month-end coefficients fixed.

\textbf{Residual rank target.} Gu, Kelly, and Xiu (2020) predict excess returns, which load on high-beta tails. Extending the residual-return principle of Blitz, Huij, and Martens (2011), the label is the centred cross-sectional percentile rank of the CAPM beta-adjusted return:
\begin{align}
\tilde{r}_{i,t+1} &= r_{i,t+1}^e - \beta_{i,t}^{60\text{m}}\, r_{m,t+1}^e, \nonumber\\
y_{i,t+1} &= 100\left(\mathrm{PctRank}_t(\tilde{r}_{i,t+1}) - \tfrac{1}{2}\right). \label{eq:resrank}
\end{align}

\textbf{Value-weighted Encompassing LASSO.} Building on Rapach, Strauss, and Zhou (2010) and Han et al.\ (2019):
\begin{enumerate}
\item \textit{Weights and penalty selection.} Each monthly regression weights stock $i$ in proportion to market equity, rescaled to mean one, so microcaps cannot dominate. Coordinate descent (Friedman, Hastie, and Tibshirani, 2010) traces a $\lambda$ path and selects $\lambda_\tau^\star$ by the corrected AIC of Hurvich and Tsai (1989).
\item \textit{Two branches.} S-LASSO fits all predictors jointly and averages coefficients over the expanding history, $\bar{\bm{\beta}}_t^{\text{S}} = \frac{1}{t-1}\sum_{\tau<t}\hat{\bm{\beta}}_\tau^{\text{S}}$. C-LASSO forms out-of-sample univariate forecasts $u_{i,j,\tau}$, selects the set $\mathcal{S}_\tau$ with a non-negative LASSO on those forecasts, and equal-weights the survivors: $\hat{y}_{i,t}^{\text{C}} = |\mathcal{S}_{t-1}|^{-1}\sum_{j \in \mathcal{S}_{t-1}} u_{i,j,t}$.
\item \textit{Forecast encompassing.} Following Chong and Hendry (1986), let $e_{i,\tau}^{\text{S}}$ be the S-LASSO error and $d_{i,\tau} = \hat{y}_{i,\tau}^{\text{C}} - \hat{y}_{i,\tau}^{\text{S}}$. With value-weighted demeaned series $\tilde{d}$ and $\tilde{e}^{\text{S}}$,
\begin{align}
\hat{\lambda}_{\tau} &= \mathrm{clip}\!\left(\frac{\sum_i v_{i,\tau}\tilde{d}_{i,\tau}\tilde{e}_{i,\tau}^{\text{S}}}{\sum_i v_{i,\tau}\tilde{d}_{i,\tau}^2},\, 0,\, 1\right), \nonumber\\
\hat{y}_{i,t} &= (1 - \bar{\lambda}_t)\,\hat{y}_{i,t}^{\text{S}} + \bar{\lambda}_t\,\hat{y}_{i,t}^{\text{C}}, \label{eq:elasso}
\end{align}
where $\bar{\lambda}_t$ is the expanding mean of $\hat{\lambda}_{\tau}$.
\end{enumerate}

\textbf{Market-state dial and horizon scaling.} Extending Pollet and Wilson (2010), each formation falls in one cell of a $3 \times 2$ grid: terciles of the 63-day correlation share (equal-weighted index variance over mean stock variance, with cuts from past data only) crossed with the sign of the trailing 12-month HML return. Using credibility shrinkage in the spirit of B\"uhlmann and Straub (1970) with prior weight $N_0 = 12$, the cell's past rank IC is shrunk toward the unconditional rank IC:
\begin{equation}
\rho_t = \mathrm{clip}\!\left(\frac{n_{\text{cell}}\,\overline{\mathrm{IC}}_{\text{cell},t} + N_0\,\overline{\mathrm{IC}}_{\text{all},t}}{(n_{\text{cell}} + N_0)\,\overline{\mathrm{IC}}_{\text{all},t}},\, 0,\, 1.5\right). \label{eq:rho_dial}
\end{equation}
The dial is 1 until 36 past formations exist or when the unconditional IC is at most 0.005. With median-centred score $\tilde{m}_{i,t}$, scale $s_f = 0.0484$, and decay $\kappa = 0.00217$, in the spirit of G\^arleanu and Pedersen (2013), the $H$-session forecast is
\begin{equation}
\mu_{i,t}^{(H)} = \rho_t\, s_f\, \frac{\tilde{m}_{i,t}}{100}\,\frac{1 - e^{-\kappa H}}{1 - e^{-21\kappa}}.
\end{equation}

\section{Factor and Semivariance Risk Model}

At each formation the $H$-session covariance is $\bm{\Sigma}_t^{(H)} = \mathbf{B}_t \mathbf{F}_t^{(H)}\mathbf{B}_t^\top + \mathbf{D}_t^{(H)}$, built in six steps.
\begin{enumerate}
\item \textit{Hedge beta.} Whereas Vasicek (1973) shrinks toward one cross-sectional mean, we shrink 252-day betas toward their GICS sector mean and average the result equally with the 60-month beta.
\item \textit{Factor returns.} Daily returns are regressed on one beta column, GICS sector indicators, and five rank-standardised styles (size, value, momentum, profitability, investment) by $\sqrt{\mathrm{ME}}$-weighted least squares, with a cap-weighted sector sum-to-zero constraint (Menchero, Orr, and Wang, 2011).
\item \textit{Horizon covariance.} With an 84-day half-life, lag-$j$ autocovariances $\bm{\Gamma}_j$ extend Newey and West (1987) to the holding horizon through
\begin{equation}
\widetilde{\mathbf{F}}^{(H)} = H\bm{\Gamma}_0 + \sum_{j=1}^{H-1}(H - j)\big(\bm{\Gamma}_j + \bm{\Gamma}_j^\top\big),
\end{equation}
followed by an eigenvalue floor at $10^{-8}\lambda_{\max}$ and 20\% diagonal shrinkage.
\item \textit{Specific risk.} On 2017--2020 formations a Patton (2011) QLIKE race with Diebold and Mariano (1995) statistics selects between HAR-RV (Corsi, 2009) and EWMA forecasts (half-lives 84 and 168 days), with 25\% shrinkage toward the sector median and a floor at the 5th percentile. We also estimate separate downside and upside residual semivariances, $d_{i,t}^-$ and $d_{i,t}^+$.
\item \textit{Volatility regime and residual factors.} A 42-day half-life volatility regime adjustment rescales factor and specific variances. Following Connor and Korajczyk (1986), five principal components of 252-day residuals are appended to $\mathbf{B}_t$.
\item \textit{Stress mixture.} A two-state Markov chain in the spirit of Hamilton (1989), on 21-day realised market variance, gives a stress probability $p_t$ and $\mathbf{F}_t \leftarrow (1 - p_t)\mathbf{F}_t + p_t\mathbf{F}_{\text{stress},t}$. The stressed matrix moves correlations halfway to $\pm 1$ and scales factor volatilities by 1.5.
\end{enumerate}

\section{Portfolio Construction and Execution}

\textbf{Signed QP.} With $\mathbf{q} = \mathbf{w} - \mathbf{w}_{\text{pre}}$, the final stage solves
\begin{align}
\min_{\mathbf{w}} \;\; & -\bm{\mu}_t^{(H)\top}\mathbf{w} + \frac{\gamma}{2}\Big(\mathbf{w}^\top \mathbf{B}_t \mathbf{F}_t^{(H)}\mathbf{B}_t^\top \mathbf{w} + H\sum_i d_{i,t}^{\pm} w_i^2\Big) \nonumber\\
& + \lambda_\beta\big(\bm{\beta}_{\text{hedge}}^\top \mathbf{w}\big)^2 + \frac{H}{252}\mathbf{b}_t^\top \mathbf{w}^- + \sum_i \frac{\lambda_u H d_{i,t}}{c_{i,t}}w_i^2 \nonumber\\
& + \mathrm{TC}(\mathbf{q}) + M_1\|\bm{\xi}_1\|_1 + M_2\|\bm{\xi}_2\|_1, \label{eq:fiam_qp}
\end{align}
with $\gamma = 0.25$, $\lambda_\beta = 10$, feature-coverage penalty $\lambda_u = 0.25$ ($c_{i,t} \in [0.1, 1]$ is the share of non-missing characteristics), and exact slack penalties $M_1 = 5$ and $M_2 = 1$. Here $d_{i,t}^{\pm}$ is the downside semivariance for longs and the upside semivariance for shorts, floored at half the symmetric variance. An earlier free-sign stage uses $\gamma = 1$. To handle the 500-name limit and a 10\,bp minimum position, we solve the continuous relaxation over eligible candidates, keep at most 440 names, lock their signs, and re-solve. Sixty names of headroom are reserved for the overlap of exits and entries. All QPs are solved with Clarabel (Goulart and Chen, 2024) through CVXPY (Diamond and Boyd, 2016).

\textbf{Constraints.} The core problem imposes gross exposure at most 180\%; sector net within $\pm 1\%$; style bands of $0.01/\sigma_{\text{factor},k}$; ex-ante volatility at most 12\% and stressed volatility at most 24\%; long caps of 4.25\% and short caps of 1\%, each also limited by five days at 10\% of average daily volume; and beta limits on three estimators: $|\bm{\beta}_{\text{hedge}}^\top \mathbf{w}| \le 0.02$, $|\bm{\beta}_{60\text{m}}^\top \mathbf{w}| \le 0.05$, and $|\bm{\beta}_{\text{WD}}^\top \mathbf{w}| \le 0.05$, where $\bm{\beta}_{\text{WD}}$ averages the Welch (2022) and Dimson (1979) betas.

\textbf{Costs and borrow.} Half-spreads scale each stock's 21-day high--low estimate (Corwin and Schultz, 2012) by a ratio calibrated on 2015--2020 TAQ effective spreads within size groups. Executions pay brokerage and regulatory fees, the half-spread, and square-root impact (Almgren et al., 2005), which the QP approximates with a matched quadratic term. Short borrow is tiered on short interest over shares (D'Avolio, 2002): 30\,bps a year below 5\%, then 1\%, 3\%, and 8\% at the 5\%, 10\%, and 20\% thresholds.

\textbf{Submitted book and ladder.} The submitted book scales core targets by $1.08$, with formation gross capped at 188\% and caps of 3.9\% long and 1.08\% short. Trades run in equal fifths over five sessions, with a 10\% daily ADV limit. Guards trim gross from 190\% to 185\% and any long from 4.4\% to 3.9\%, and cover any short whose capitalisation falls below a \$550m buffer. Eight local language-model agents read filings, draft commentary, measure crowding, and challenge the book; none sets a weight.

\begin{table*}[t]
\centering
\caption{Headline results, 4 January 2021 to 31 August 2026 (68 months, 1,421 sessions), after modelled costs. Volatility, Sharpe, Sortino, beta, and alpha use monthly returns; Sortino uses a 0\% target. S\&P 500 figures use the price index. Source: locked results in the public repository.}
\label{tab:results}
\footnotesize
\setlength{\tabcolsep}{5pt}
\begin{tabular}{lrrr}
\toprule
Metric & QUARCC Purple & T-bill $+4\%$ & S\&P 500 \\
\midrule
Ending value of \$1.000bn & \$1.467bn & \$1.509bn & \$2.046bn \\
CAGR & 7.00\% & 7.53\% & 13.47\% \\
Volatility & 3.96\% & --- & 15.10\% \\
Sharpe ratio & 0.91 & --- & 0.70 \\
Sortino ratio & 4.21 & --- & 1.51 \\
Maximum drawdown & $-2.99\%$ & --- & $-25.43\%$ \\
Ulcer index & 0.52\% & --- & 8.10\% \\
Beta to the S\&P 500 (s.e.) & 0.012 (0.030) & --- & 1.00 \\
Alpha to the S\&P 500 ($t$) & 3.47\% a year (2.45) & --- & --- \\
Information ratio vs.\ T-bill $+4\%$ (s.e.) & $-0.11$ (0.42) & --- & --- \\
Worst month / worst day & $-1.43\%$ / $-1.37\%$ & --- & --- \\
Months positive & 48 of 68 & --- & --- \\
\bottomrule
\end{tabular}
\qquad
\begin{tabular}{lrr}
\toprule
Year & Fund & T-bill $+4\%$ \\
\midrule
2021 & 6.98\% & 4.12\% \\
2022 & 1.34\% & 6.19\% \\
2023 & 15.10\% & 9.45\% \\
2024 & 11.16\% & 9.34\% \\
2025 & 2.08\% & 8.37\% \\
2026 (to Aug.) & 3.62\% & 5.21\% \\
\bottomrule
\end{tabular}
\end{table*}

\section{Results}

Table~\ref{tab:results} reports results recomputed from the repository's locked daily ledger. The rule checks record zero breach days over the window, with 246--357 names held and gross exposure averaging 136.5\% (maximum 192.9\%, final 186.0\%). Single-name exposure reached at most 4.6\% and net exposure stayed between $-3.1\%$ and $+2.2\%$.

Of the \$467.07m net profit, stock positions contributed \$257.7m, interest on cash \$239.5m, and costs took \$30.1m: \$9.5m of trading (14\,bps a year) and \$20.6m of borrow (29\,bps a year). Interest is thus nearly half of gross profit, which is why the T-bill-based hurdle is the appropriate comparison. In 2022, when the S\&P 500 total return was $-18.18\%$, the fund returned $+1.34\%$.

\section{Limitations}

\textit{The fund misses its benchmark.} The 7.00\% CAGR trails the 7.53\% hurdle by 53\,bps a year. Costs equal 72\% of the \$41.6m cumulative shortfall, and the information ratio before costs is about zero. The standard error of 0.42 means the data cannot distinguish a modest edge from none.

\textit{It is not a clean holdout.} Forecasts are walk-forward, but the construction settings, including the final gross scale, were chosen with 2021--2026 visible. Sharpe ratios and drawdowns should be read with that selection in mind.

\textit{It lags in rallies.} The fund trailed its hurdle in 2022, 2025, and 2026 to date. In 2025 it returned 2.08\% against a hurdle of 8.37\%. Gross exposure averaged 136.5\%, well below the 200\% limit, because the optimiser charges entry costs against a horizon-scaled forecast and builds slowly.

\textit{Costs are modelled.} No order was executed. Fills, borrow availability, and impact at \$1 billion may differ from the model.

\textit{The live tracker is a paper mark.} QUARCC's public paper-fund page marks a single frozen 24 August 2026 formation forward with market data. It does not execute trades and is not a live track record.

\section*{Code and Data}

Code is available under the MIT License at \href{https://github.com/NicholasXydis/QuarccPurple}{github.com/NicholasXydis/QuarccPurple}, with the fact sheet, the competition deck, and the locked results. Rule checks and headline recomputation run without licensed data. Licensed inputs from WRDS and the organisers are required for a full rebuild and cannot be shared.

\vspace{3pt}
\section*{References}
\scriptsize
\begin{list}{}{\setlength{\leftmargin}{1.3em}\setlength{\itemindent}{-1.3em}\setlength{\itemsep}{1.0pt}\setlength{\parsep}{0pt}}
\item Almgren, Robert, Chee Thum, Emmanuel Hauptmann, and Hong Li, 2005, Direct estimation of equity market impact, \textit{Risk} 18, 58--62.
\item Bernard, Victor L., and Jacob K. Thomas, 1989, Post-earnings-announcement drift: Delayed price response or risk premium?, \textit{Journal of Accounting Research} 27, 1--36.
\item Blitz, David, Joop Huij, and Martin Martens, 2011, Residual momentum, \textit{Journal of Empirical Finance} 18, 506--521.
\item Boehmer, Ekkehart, Zsuzsa R. Huszar, and Bradford D. Jordan, 2010, The good news in short interest, \textit{Journal of Financial Economics} 96, 80--97.
\item B\"uhlmann, Hans, and Erwin Straub, 1970, Glaubw\"urdigkeit f\"ur Schadens\"atze, \textit{Mitteilungen der Vereinigung Schweizerischer Versicherungsmathematiker} 70, 111--133.
\item Chong, Yock Y., and David F. Hendry, 1986, Econometric evaluation of linear macro-economic models, \textit{Review of Economic Studies} 53, 671--690.
\item Connor, Gregory, and Robert A. Korajczyk, 1986, Performance measurement with the arbitrage pricing theory, \textit{Journal of Financial Economics} 15, 373--394.
\item Corsi, Fulvio, 2009, A simple approximate long-memory model of realized volatility, \textit{Journal of Financial Econometrics} 7, 174--196.
\item Corwin, Shane A., and Paul Schultz, 2012, A simple way to estimate bid-ask spreads from daily high and low prices, \textit{Journal of Finance} 67, 719--760.
\item D'Avolio, Gene, 2002, The market for borrowing stock, \textit{Journal of Financial Economics} 66, 271--306.
\item Diamond, Steven, and Stephen Boyd, 2016, CVXPY: A Python-embedded modeling language for convex optimization, \textit{Journal of Machine Learning Research} 17, 1--5.
\item Diebold, Francis X., and Roberto S. Mariano, 1995, Comparing predictive accuracy, \textit{Journal of Business \& Economic Statistics} 13, 253--263.
\item Dimson, Elroy, 1979, Risk measurement when shares are subject to infrequent trading, \textit{Journal of Financial Economics} 7, 197--226.
\item Friedman, Jerome, Trevor Hastie, and Rob Tibshirani, 2010, Regularization paths for generalized linear models via coordinate descent, \textit{Journal of Statistical Software} 33, 1--22.
\item G\^arleanu, Nicolae, and Lasse Heje Pedersen, 2013, Dynamic trading with predictable returns and transaction costs, \textit{Journal of Finance} 68, 2309--2340.
\item Goulart, Paul J., and Yuwen Chen, 2024, Clarabel: An interior-point solver for conic programs with quadratic objectives, arXiv preprint arXiv:2405.12762.
\item Green, Jeremiah, John R. M. Hand, and X. Frank Zhang, 2017, The characteristics that provide independent information about average U.S.\ monthly stock returns, \textit{Review of Financial Studies} 30, 4389--4436.
\item Gu, Shihao, Bryan Kelly, and Dacheng Xiu, 2020, Empirical asset pricing via machine learning, \textit{Review of Financial Studies} 33, 2223--2273.
\item Hamilton, James D., 1989, A new approach to the economic analysis of nonstationary time series and the business cycle, \textit{Econometrica} 57, 357--384.
\item Han, Yufeng, Ai He, David E. Rapach, and Guofu Zhou, 2019, Firm characteristics and expected stock returns, Working paper.
\item Hurvich, Clifford M., and Chih-Ling Tsai, 1989, Regression and time series model selection in small samples, \textit{Biometrika} 76, 297--307.
\item Jensen, Theis I., Bryan Kelly, and Lasse Heje Pedersen, 2023, Is there a replication crisis in finance?, \textit{Journal of Finance} 78, 2465--2518.
\item Menchero, Jose, D. J. Orr, and Jun Wang, 2011, The Barra US Equity Model (USE4) methodology notes, MSCI Barra.
\item Newey, Whitney K., and Kenneth D. West, 1987, A simple, positive semi-definite, heteroskedasticity and autocorrelation consistent covariance matrix, \textit{Econometrica} 55, 703--708.
\item Patton, Andrew J., 2011, Volatility forecast comparison using imperfect volatility proxies, \textit{Journal of Econometrics} 160, 246--256.
\item Pollet, Joshua M., and Mungo Wilson, 2010, Average correlation and stock market returns, \textit{Journal of Financial Economics} 96, 364--380.
\item Rapach, David E., Jack K. Strauss, and Guofu Zhou, 2010, Out-of-sample equity premium prediction: Combination forecasts and links to the real economy, \textit{Review of Financial Studies} 23, 821--862.
\item Stambaugh, Robert F., Jianfeng Yu, and Yu Yuan, 2012, The short of it: Investor sentiment and anomalies, \textit{Journal of Financial Economics} 104, 288--302.
\item Vasicek, Oldrich A., 1973, A note on using cross-sectional information in Bayesian estimation of security betas, \textit{Journal of Finance} 28, 1233--1239.
\item Welch, Ivo, 2022, Simply better market betas, \textit{Critical Finance Review} 11, 37--64.
\end{list}

\end{document}
